\clearpage
\VerdierPhysicalPageAnchor{0218}%
\begingroup\small\setlength{\parskip}{0pt}\setlength{\emergencystretch}{1.5em}
\setlength{\abovedisplayskip}{4pt plus 2pt minus 2pt}%
\setlength{\belowdisplayskip}{4pt plus 2pt minus 2pt}%

\VerdierSemanticAnchor{ch3d-3-3-1-2}
\noindent\textit{$S_i^1*f_{i-1}$，$1<i\leq n$，
$f_i*{S''}_i^1$，$1\leq i<n$，${S'}_i^1*g_{i-1}$，$1<i\leq n$，
$g_i*{S''}_i^1$，$1\leq i<n$，均有定义（(3.2.8.1) 与 (3.2.8.2)），
并且满足：}

\begin{equation}
\left\{
\begin{aligned}
\delta^1({S''}_n^1)
  &=\delta^1(S_n^1*f_{n-1}),\\
\delta^1({S''}_n^1)
  &=\delta^1({S'}_n^1*g_{n-1}),\\
\delta^1(f_i*{S''}_i^1)
  &=\delta^1(S_i^1*f_{i-1}),&&1<i\leq n-1,\\
\delta^1(g_i*{S''}_i^1)
  &=\delta^1({S'}_i^1*g_{i-1}),&&1<i\leq n-1,\\
\delta^1(f_1*{S''}_1^1)
  &=\delta^1(S_1^1),\\
\delta^1(g_1*{S''}_1^1)
  &=\delta^1({S'}_1^1).
\end{aligned}
\right.
\tag{3.3.1.2}
\end{equation}

\VerdierSemanticAnchor{ch3d-3-3-1-proof}
事实上，断言 a) 由 (3.2.2) 与 (3.2.5) 立即得到，而断言 b) 由映射
$E_{\mathcal A}^{n}(X,Y)\longrightarrow
\operatorname{Ext}_{\mathcal A}^{n}(X,Y)$ 所诱导的等价关系之描述得到
（(3.2.2) 与 (3.2.9)）。

\VerdierSemanticAnchor{ch3d-3-3-2}
\medskip\noindent\textbf{3.3.2.}\quad
回忆 (1.2.11)，范畴 $\mathcal A^{\operatorname{Ext}}$ 是
$D(\mathcal A)$ 的满子范畴，它由 $\mathcal A$ 的对象的各个平移定义。
设 $F\colon\mathcal A^{\operatorname{Ext}}\longrightarrow\mathcal C$ 是一个
函子（取值于一个不必为加法范畴的范畴）。对每个整数 $n$，以
$F^n\colon\mathcal A\longrightarrow\mathcal C$ 表示函子
$X\longmapsto F(X[n])$。同样，对 $\mathcal A$ 的每条正合列
$S^1=0\longrightarrow X\longrightarrow Y\longrightarrow Z
\longrightarrow0$，以
$\delta_F^n(S^1)\colon F^n(Z)\longrightarrow F^{n+1}(X)$ 表示态射

\VerdierSemanticAnchor{ch3d-3-3-2-delta-definition}
\[
F\bigl(\delta^1(S^1)[n]\bigr)
\colon F(Z[n])\longrightarrow F(X[n+1]).
\]

\VerdierSemanticAnchor{ch3d-3-3-2-sfl}
函子列 $F^n$，$n\in\mathbb Z$，连同态射 $\delta_F^\ast$，具有下列性质
(3.2.9)：

\noindent\textbf{SFL I)} 每个交换图：

\VerdierSemanticAnchor{ch3d-3-3-2-sfl-i-source-diagram}
\begin{center}
\resizebox{0.86\textwidth}{!}{%
$\begin{tikzcd}[ampersand replacement=\&,column sep=3.0em,row sep=2.8em]
S=0 \arrow[r]
  \& X \arrow[r] \arrow[d,"f"']
  \& Y \arrow[r] \arrow[d]
  \& Z \arrow[r] \arrow[d,"g"]
  \& 0 \\
{S'}=0 \arrow[r]
  \& X' \arrow[r]
  \& Y' \arrow[r]
  \& Z' \arrow[r]
  \& 0
\end{tikzcd}$%
}
\end{center}
\endgroup

